% =====================================================================
%  Chapter 25 — Why the Answer Is a Portfolio
% =====================================================================

\chapter{Why the Answer Is a Portfolio}
\label{ch:portfolio}

The last three chapters of this book are about decisions. This chapter gives
the frame for them: why the strategy that the evidence implies is a
portfolio of differentiated experiments rather than a race, and what follows
from that for how the field should allocate effort.

\section{Four arguments for the portfolio}Each of the following is an
independent argument for the same structure, and their conjunction is the
strongest strategic claim in this book.

The first is \emph{geometric}. The exclusion curves of
Chapter~\ref{ch:direct} show that no experiment covers the remaining
parameter space. The curves have different shapes because the experiments
respond to different couplings, and their complementarity is a fact about
the experiments rather than a choice. The Particle Data Group review states
the consequence directly: the long-term ceiling of the current generation is
of order $3\times 10^{-49}\ \cmtwo$, and the parameter space beyond it must
be attacked by other means \idx{Particle Data Group}.

The second is \emph{strategic}. The 2023 Particle Physics Opportunity
Assessment states that no single experiment can probe the remaining natural
parameter space and that a portfolio of complementary approaches is required
\idx{APPEC}. The 2022 Snowmass white papers make the same argument from the
theory side, identifying the regions where different experiments overlap and
the regions where only one reaches \idx{Snowmass complementarity}. A
portfolio is not a compromise between specialists; it is the structure the
parameter space has.

The third is \emph{epistemic}. The anomaly record of Chapter~\ref{ch:indirect}
shows that the binding constraint on the field's progress is replication, not
sensitivity. Four significant claims in five years did not replicate, and
each of them was made by a single group on a single analysis. The
consequence is that the field's scarcest resource is not sensitivity but
independent replication, and a portfolio of independent experiments
produces replication as a by-product.

The fourth is \emph{epistemic in the other direction}. The origins question
of Chapter~\ref{ch:origin} cannot be settled by any experiment. It is
settled by a combination: the particle's properties, the cosmological
history, and the elimination of inconsistent combinations. Each of those
requires a different instrument, and a portfolio is the only structure that
contains all three.

\section{What the portfolio is not}It is worth being explicit about three
misunderstandings, because each of them has been proposed.

The portfolio is not a way of avoiding a decision. Every candidate in
Chapter~\ref{ch:criteria} has a decision criterion attached, and each
experiment in the roadmap has a kill condition. The point of a portfolio is
that the decision criteria are applied to each branch independently, not that
the decision is postponed.

It is not an argument for more of the same. The dominant experiment type in
the field --- the deep-underground liquid-xenon nuclear-recoil search ---
has reached the point where a ten-fold increase in exposure buys a factor of
a few in cross section. The portfolio's diversification is into targets
(light nuclei), into recoils (electrons), into directions, into devices
(qubits), and into the sky and the early universe.

And it is not a reason to keep everything alive. A portfolio with no kill
conditions is a research programme, not a strategy. The criteria in
Chapter~\ref{ch:criteria} are what distinguish the two.

\section{How to divide the effort}A reasonable division, on the logic of the
preceding chapters, is by the bottleneck each technique addresses rather than
by the candidate each technique favours.

\begin{itemize}
\item The \emph{deep underground} programme has reached its floor in the
nuclear-recoil channel and should be reoriented towards the light-target,
electron-recoil and low-mass regimes where the floor has not been reached.
\item The \emph{astrophysical} programme --- dwarf galaxies, clusters, the
Milky Way's centre, the halo's velocity distribution --- is where the
discrimination between candidate classes is highest and where the
systematics are improving with survey quality.
\item The \emph{early universe} programme --- the primordial spectrum, the
CMB spectral distortions, the thermal history --- is the cheapest route to
the origins question and the most neglected relative to its leverage.
\item The \emph{device} programme --- quantum sensors, resonant conversion,
directional detectors --- is where the qualitatively new measurements are,
and where the instrumentation culture is weakest.
\end{itemize}

\section{The asymmetry of the field's resources}It is worth saying plainly
what the portfolio requires that the field does not currently have.

Direct detection has consumed the resources: the tonne-scale experiments
cost hundreds of millions of dollars each and the field's growth has tracked
them. The astrophysical programme is comparatively cheap and has been
chronically under-resourced relative to the number of questions it can
answer. The early-universe programme requires no new facility at all, only
data and analysis, and it is the smallest community in the book.

The strategy that the evidence implies is therefore not primarily a
strategy of building more large detectors. It is a strategy of using the
existing sky and the existing universe better, while building the smaller
and more diverse instruments that the floor of the large ones demands.

\begin{keytakeaways}
\item Four independent arguments --- geometric, strategic, epistemic and
origins-driven --- all require a portfolio of differentiated experiments.
\item A portfolio is not a way of avoiding decisions; it is a structure in
which each branch has its own criteria and its own kill conditions.
\item The effort should be divided by bottleneck, not by favoured candidate.
\item The field's resources are concentrated in the channel that has reached
its floor, and the cheapest high-leverage programme is the one that uses
existing data.
\end{keytakeaways}
